% Generated draft: requires mathematical and source review.
\documentclass[11pt]{article}
\usepackage[margin=25mm]{geometry}
\usepackage{fontspec}
\setmainfont{Latin Modern Roman}
\usepackage{amsmath,amssymb,mathtools,unicode-math}
\setmathfont{Latin Modern Math}
\usepackage{enumitem,needspace,xurl,hyperref}
\setlength{\parindent}{0pt}
\setlength{\parskip}{5pt}
\setlength{\emergencystretch}{3em}
\newcommand{\N}{\mathbb{N}}
\newcommand{\Z}{\mathbb{Z}}
\newcommand{\Q}{\mathbb{Q}}
\newcommand{\R}{\mathbb{R}}
\newcommand{\C}{\mathbb{C}}
\newcommand{\F}{\mathbb{F}}
\newcommand{\A}{\mathbb{A}}
\newcommand{\PP}{\mathbb{P}}
\newcommand{\E}{\mathbb{E}}
\newcommand{\eps}{\varepsilon}
\newcommand{\dd}{\,\mathrm{d}}
\begin{document}
\section*{UnsolvedMath ID: 6700017}
\textbf{Source problem code:} catalogue entry
\subsection{Definitions and mathematical statement}
The tagged problem is Conjecture 6, ``ISC: Singularities are Irrelevant.'' It asks whether the Schoen--Yau minimal-subvariety descent method extends to singular minimal subvarieties. In this context descent means passing to lower-dimensional minimal objects while retaining the geometric/topological information needed for scalar-curvature arguments. The record gives no precise class of singular subvarieties, admissible modifications, dimension range, or exact preserved conclusion. It does not assert that positive scalar curvature makes all minimal-hypersurface singularities disappear. A fully quantified ISC theorem cannot be recovered from the slogan alone.

The supplied text then runs into the next section, on flatly twisted spinors over tori. This is adjacent context, not an additional explicit formulation of ISC. It describes a flat even-dimensional torus, whose positive and negative Dirac kernels consist of parallel spinors of equal dimension, so the ordinary index vanishes. Its displayed kernel dimension appears as the damaged string ``2n−1''; no exponent or corrected dimension is inferred.

For the usual closed spin even-dimensional setting, the adjacent construction considers characters $\tau:\pi_1(X)\to U(1)$, the associated flat line bundles $l_\tau$, and twisted operators
\[
 D_\tau^\pm:C^\infty(S^\pm\otimes l_\tau)
                \longrightarrow C^\infty(S^\mp\otimes l_\tau).
\]
It places their family index in $K^0(T^m)$, with $m=\operatorname{rank}_{\Q}H^1(X;\Q)$. A full character space can have finite components; the identity component is a torus, so the source's identification of the entire space with $T^m$ needs a torsion convention. Individual kernel dimensions can jump, so the virtual family index does not require the two kernels to be ordinary vector bundles of constant rank as the informal source wording suggests.

The adjacent passage calls $X$ over-torical when it has a nonzero-degree map to $T^n$, and reports a nonzero family index in this case. A standard cohomological expression of that degree condition, for a closed oriented $n$-manifold, is existence of integral degree-one classes $a_1,\ldots,a_n$ with
\[
 \langle a_1\smile\cdots\smile a_n,[X]\rangle\ne0.
\]
The extracted text instead speaks of ``homology classes'' and writes an ambiguous $H1$ with intersection notation; its grading must be checked. These adjacent family-index claims are retained as source context with the required compactness/spin conventions made explicit for review, not as a theorem proved here or a replacement for the missing precise ISC statement.
\subsection{Short English statement}
Can scalar-curvature descent through minimal subvarieties be carried out despite singularities? The source's ISC statement is programmatic and lacks an exact implication. The extraction also includes the next section's torus family-index discussion, which is kept distinct from the tagged question.

The source reports Schoen--Yau announcements roughly 35 years earlier, Lohkamp's proposed approach, and a recent alternative proof of a version by Schoen--Yau. Gromov explicitly adds that he has not studied the Lohkamp and Schoen--Yau papers in depth. The supplied summary describes low-dimensional success and higher-dimensional partial progress, associates the issue with aspherical positive-scalar-curvature questions, and calls the general singularity-descent strategy unresolved in the 2020s, especially from dimension eight. Those are supplied assessments; the incomplete author PDF's current status needs review.
\subsection{Sources}
\begin{itemize}
\item Local UnsolvedMath registry, permanent ID 6700017, source code catalogue entry, scalar-curvature Question [?17], proposed by Misha Gromov, 2017. Dataset curation and metadata: Ulam AI and UnsolvedMath Contributors, CC BY 4.0.
\item Misha Gromov, \emph{101 Questions, Problems and Conjectures around Scalar Curvature} (2017), printed tag [?17], occurrence 17, PDF page 21, Conjecture 6 and extracted adjacent Section 7, \url{https://www.ihes.fr/~gromov/wp-content/uploads/2018/08/101-problemsOct1-2017.pdf}; normalized pdf-text extraction, accessed 2026-07-29.
\item Source references [110], [114], Schoen--Yau announcements; [115], an alternative proof of a version of irrelevance. Full entries are absent.
\item Source reference [86], Lohkamp, \emph{The Higher Dimensional Positive Mass Theorem II} (2016), with earlier papers referenced there, as supplied. No further publication details are inferred.
\item The adjacent family-index passage attributes a 1972 result to ``Lustig'' and ends with a truncated reference marker. The name and citation are not corrected or completed without source evidence.
\item Source-specific provenance: the author-hosted PDF is incomplete and publicly available; the numbered passage is an open-item tag, while current status and statement-reuse terms NEEDS\_REVIEW. These qualifications are separate from the dataset licence.
\end{itemize}
\end{document}
